\documentclass[11pt]{article}
\usepackage{amsmath, amssymb, amsthm}
\usepackage{geometry}
\usepackage{hyperref}
\geometry{letterpaper, margin=1in}

\title{Algebraic Characterization of Safe Hook Errors in ZX-Calculus}
\author{}
\date{\today}

\begin{document}

\maketitle

\section{Definition of Safe Hook Errors}
In fault-tolerant (FT) state preparation circuits synthesized via ZX-calculus, diagrammatic rewrites strictly optimize the topology of a tensor network. To extract the graph into a compilable circuit, high-degree spiders (representing large stabilizer generators) must be split into lower-degree components.

Splitting a spider creates an internal edge. A single physical fault occurring on this internal edge propagates to the boundary data block as a multi-qubit Pauli error, denoted $E$. We operate under a fault-tolerant distance $d$. Because the diagram prepares specific logical states, it is stabilized by a preparation group $\mathcal{S}_{\text{prep}}$.

An error $E$ resulting from a split is defined as \textbf{strictly safe} if it acts identically to a single physical fault (weight $\le 1$) on the target state. Mathematically, the equivalent physical weight $w_{\text{min}}^*$ must satisfy:
\begin{equation}
    w_{\text{min}}^*(E) = \min_{S \in \mathcal{S}_{\text{prep}}} \text{weight}(E \oplus S) \le 1
\end{equation}
If this holds, the hook error is harmless because it is algebraically identical (modulo the state symmetries) to a single physical fault. By the triangle inequality, combinations of strictly safe hook errors cannot unexpectedly amplify to cause logical failures.

\section{Intuition for Checking Safety}
The algebraic requirement $w_{\text{min}}^*(E) \le 1$ implies that there exists some stabilizer $S \in \mathcal{S}_{\text{prep}}$ and a physical error $e$ of weight $\le 1$ such that:
\begin{equation}
    E = S \oplus e
\end{equation}

Since $E$ is generated by splitting a specific generator with support $G$, the propagated error $E$ is strictly localized to $G$. Therefore, on the complement set $I = G^C$, the error $E$ has zero support ($E_I = 0$).

This localization provides a powerful algebraic constraint:
\begin{equation}
    0 = S_I \oplus e_I \implies S_I = e_I
\end{equation}
Checking if a split is safe reduces to searching for a stabilizer $S$ that matches a weight $\le 1$ error strictly on the outside qubits $I$. 

\section{Efficient Description of All Safe Hook Errors}
Naively checking all possible subsets of $G$ requires evaluating $2^{|G|}$ configurations, which scales exponentially with the maximum stabilizer weight. However, we can efficiently describe the exact space of safe splits in strictly polynomial $\mathcal{O}(N^3)$ time by inverting the problem over GF(2).

Let $M_{\text{prep}}$ be the binary generator matrix of $\mathcal{S}_{\text{prep}}$, and $M_I$ be its restriction to the columns in $I$. Because a stabilizer $S$ is a linear combination of rows $x^T M_{\text{prep}}$, the condition $S_I = e_I$ becomes the linear system:
\begin{equation}
    x^T M_I = e_I
\end{equation}

The solutions to this system geometrically form a union of affine subspaces (cosets):
\begin{enumerate}
    \item \textbf{The Base Space ($V_0$):} If $e_I = 0$, $x$ belongs to the left null-space of $M_I$. The projection of these stabilizers onto $G$ defines a vector space $V_0 = \{ x^T M_G \mid x^T M_I = 0 \}$. These correspond to splits that are natively safe without needing to pair with any outside faults.
    \item \textbf{The Shift Vectors ($U$):} For each possible solvable weight-1 fault outside $G$ (where $e_I$ has weight 1), solving the linear system yields a particular base solution $u_q$.
\end{enumerate}

The exact set of safe splits is strictly the affine union:
\begin{equation}
    \text{Safe Splits} = \bigcup_{u \in U} (u \oplus V_0)
\end{equation}
This $(V_0, U)$ pair is the most compact mathematical description possible, bypassing the exponential search space entirely.

\section{Universal vs. Partial Safety}
When extracting circuits, spiders are recursively fractured into partition shapes (e.g., separating a degree-6 spider into two pieces of size 3, denoted as a $[3,3]$ partition). We classify the topological safety of these partitions into two categories:

\subsection{Universal Safety}
A partition shape is \textbf{universally safe} if \textit{any} arbitrary choice of qubits matching the chunk sizes forms a valid safe split sequence. 

Instead of generating splits, we mathematically compute the exact weight distribution of the affine space $\bigcup_{u \in U} (u \oplus V_0)$. If the number of safe splits of a specific Hamming weight $k$ perfectly equals the binomial coefficient $\binom{|G|}{k}$, then \textit{every} subset of size $k$ is guaranteed to be safe. If all prefix sums of a partition shape consist of universally safe sizes, the shape is universally safe.

\subsection{Partial Safety}
If a partition shape is not universally safe, it may still be \textbf{partially safe}. This occurs when there exists at least one specific nested sequence of safe splits matching the shape, but not all choices of qubits work. 

To quickly generate a partially safe set, we explicitly generate the members of the affine space $\bigcup (u \oplus V_0)$. Because the dimension of $V_0$ corresponds to the number of independent stabilizers fully supported on $G$ (which is practically $O(1)$), generating this exact set is blazingly fast. We then perform a topological Depth-First Search (DFS) over the partially ordered set (poset) of these specific safe splits to extract genuine nested structural chains.

\end{document}
